\documentclass[11pt]{article}

\usepackage[margin=1in]{geometry}
\usepackage{booktabs}
\usepackage{graphicx}
\usepackage{hyperref}
\usepackage{amsmath}
\usepackage{amssymb}
\usepackage{microtype}

\title{GAX: Evidence-Calibrated System-One Models for Clinical Decision Making}
\author{Author list to be finalized before submission}
\date{}

\begin{document}
\maketitle

% IMPORTANT:
% Do not insert synthetic or placeholder numerical results.
% Every number in a results sentence/table must trace to a frozen evidence packet.

\begin{abstract}
GAX studies non-generative, typed clinical decision models that separate action
beliefs from information sufficiency and abstention. This working manuscript is
intentionally result-free until the frozen evaluation protocol is complete.
The final abstract will report only artifact-backed findings.
\end{abstract}

\section{Introduction}

Medical AI systems are commonly optimized to generate text, while many
operational and clinical-support workflows ultimately require bounded decisions:
select an allowed action, estimate risk, verify evidence support, or defer when
information is insufficient. GAX investigates whether a dedicated
System-One-style decision model can improve the accuracy--calibration--selectivity
trade-off for these tasks while retaining efficient inference.

The core design principle is that the highest-probability action is not itself
permission to act. GAX therefore represents action probabilities and information
sufficiency separately and evaluates selective prediction explicitly.

\paragraph{Contributions.}
Contribution language remains gated by the experiments in
\texttt{docs/research/PAPER\_PLAN.md}. Any contribution that fails its ablation
or held-out evaluation gate must be removed or reframed as a negative result.

\section{Related Work}

\subsection{System-One and Typed Decision Models}
Discuss Jev, Contrastive Language Models, Laya, decider, and restricted-logit
one-pass baselines. Distinguish public evidence from unpublished implementation
details.

\subsection{Clinical Encoders}
Discuss Clinical ModernBERT and BioClinical ModernBERT as efficient long-context
clinical encoders.

\subsection{Medical Agents and FHIR}
Discuss MedAgentBench and FHIR-AgentBench, focusing on interoperable EHR
retrieval and action selection.

\subsection{Evidence-Aware Medical Verification}
Discuss Med-PRM and related verifier/process-reward approaches.

\subsection{Uncertainty, Calibration, and Abstention}
Discuss proper scoring rules, neural calibration, selective prediction, and
MedQAbstain.

\section{Problem Formulation}

Let $s$ denote a clinical state, $\mathcal{A}=\{a_1,\ldots,a_K\}$ a finite set of
caller-permitted actions, and $E$ optional evidence. A GAX model estimates

\begin{equation}
p(a_i \mid s,E), \qquad i=1,\ldots,K,
\end{equation}

and an information-sufficiency quantity

\begin{equation}
p(\mathrm{sufficient}\mid s,E).
\end{equation}

A separate decision policy maps these quantities to either an allowed action or
abstention. Calibration parameters and action thresholds are fit without access
to held-out test labels.

\section{GAX}

\subsection{Architecture Search}
Report the winning architecture only after matched experiments across the
single-encoder, contrastive state/action, and packed one-pass candidates.

\subsection{Evidence-Calibrated Action Learning}
ECAL is a working objective under test:
\begin{equation}
\mathcal{L} =
w_a\mathcal{L}_{\mathrm{action}} +
w_e\mathcal{L}_{\mathrm{evidence}} +
w_p\mathcal{L}_{\mathrm{proper}} +
w_s\mathcal{L}_{\mathrm{selective}} +
w_c\mathcal{L}_{\mathrm{counterfactual}} +
w_r\mathcal{L}_{\mathrm{retention}}.
\end{equation}
This formulation is not a claimed contribution unless ablations support it.

\subsection{Selective Decision Policy}
Define confidence-only, learned-sufficiency, and calibrated selective baselines
at matched coverage.

\section{GAXBench}

Describe the benchmark suites for closed actions, abstention, evidence
interventions, counterfactual sensitivity, FHIR workflows, distribution shift,
and systems performance. Include data provenance and leakage controls.

\section{Experimental Setup}

\subsection{Models and Baselines}
Freeze versions before final test evaluation.

\subsection{Datasets and Splits}
Describe entity-disjoint splitting, transformation lineage, contamination checks,
and licensing constraints.

\subsection{Metrics}
Primary metrics include task accuracy where appropriate, negative log
likelihood, Brier score, calibration error with disclosed binning, area under the
risk--coverage curve, risk at fixed coverage, abstention metrics, evidence
discrimination, and systems metrics.

\subsection{Statistical Analysis}
Use paired comparisons, confidence intervals, multiple training seeds where
needed, and explicit failure/exclusion accounting.

\section{Results}

\textbf{TBD after frozen evidence-backed evaluation.}

\section{Ablations}

\textbf{TBD after frozen evidence-backed evaluation.}

\section{FHIR Agent Evaluation}

\textbf{TBD after frozen evidence-backed evaluation.}

\section{Calibration and Selective Risk}

\textbf{TBD after frozen evidence-backed evaluation.}

\section{Failure Analysis}

Report confident errors, over-abstention, evidence-following failures,
counterfactual insensitivity, irrelevant-edit instability, and FHIR
representation/retrieval errors without exposing restricted clinical data.

\section{Limitations, Ethics, and Intended Use}

GAX is research software and is not a medical device. The initial work does not
establish autonomous diagnosis, treatment, prescribing, emergency triage, or
clinical deployment safety. Constrained outputs prevent out-of-schema actions
but do not prevent clinically incorrect in-schema decisions.

\section{Reproducibility Statement}

Bind the paper to a repository tag, immutable model/data revisions, split
manifest hashes, commands, seeds, environment locks, hardware, raw metrics, and
figure/table generation code.

\section{Conclusion}

\textbf{TBD after the central hypotheses are resolved.}

\bibliographystyle{plain}
\bibliography{references}

\end{document}
